% ============================================================
% Capítulo: Trabajo Futuro (Bloque B.13)
% ============================================================
\chapter{Trabajo futuro}
\label{cap:trabajo-futuro}

Este capítulo enumera las líneas de trabajo que quedan explícitamente
fuera del alcance de esta Entrega Final -- no genéricas ni especulativas,
sino los gaps reales ya identificados en el propio repositorio a lo largo
de este documento. Cada una indica por qué quedó fuera de alcance y cómo
se abordaría.

\section{Completar la evidencia de usabilidad (SUS, $N\geq15$)}
\label{sec:tf-sus}

El Bloque~5 de evaluación empírica (\autoref{sec:res-usabilidad}) quedó
con $N=0$ de los $N\geq15$ participantes que exige la guía. Quedó fuera
de alcance de esta entrega porque su ejecución depende de una condición
externa al propio código: un despliegue público funcional donde
participantes reales puedan interactuar con el sistema bajo el protocolo
ya definido (plantilla de consentimiento informado, técnica de muestreo
por conveniencia dentro de la comunidad UTEQ, anonimización por código de
participante). Ese despliegue está en curso (trabajo de Cajas); una vez
disponible, la línea de trabajo es directa y ya está protocolizada: (i)
reclutar $\geq15$ participantes con roles representativos del dominio
(lector, bibliotecario), (ii) aplicar el instrumento SUS de 10 ítems
\citep{brooke1996sus} bajo el protocolo de
\texttt{docs/etica/consentimientos/plantilla.md}, (iii) calcular la
puntuación agregada y aplicar el criterio de \citet{baltes2022sampling}
sobre generalización de muestras pequeñas al interpretar el resultado, y
(iv) actualizar el \autoref{cap:resultados} y el
\autoref{cap:amenazas-validez} con el dato real obtenido -- ya asignada a
Panama según el plan del equipo (\texttt{OBS-08} en
\texttt{docs/observaciones/OBSERVACIONES.md}).

\section{Análisis estático de SQL dinámico}
\label{sec:tf-audit-sql}

El \autoref{cap:implementacion} señaló, al aclarar por qué
\texttt{@Query(nativeQuery=true)} con parámetros bindeados no viola la
prohibición de concatenación de cadenas de la guía, que un script de
análisis estático dedicado (\texttt{scripts/audit-sql-dynamic.sh}) para
detectar SQL dinámico construido por concatenación de entrada de usuario
todavía no existía en el repositorio en ese momento. Se verificó de
nuevo al escribir este capítulo (búsqueda directa en
\texttt{scripts/}) y ese script \textbf{sigue sin existir} -- no se
completó entre ambas verificaciones. Quedó fuera de alcance de esta
entrega porque ninguna tarea de este ciclo lo priorizó explícitamente por
sobre las auditorías manuales ya ejecutadas (\autoref{sec:res-seguridad}),
que cubren el mismo riesgo de forma directa aunque no automatizada. La
línea de trabajo futura es concreta: un script que recorra
\texttt{backend-springboot/src/main/java} buscando invocaciones a
\texttt{createNativeQuery(...)} o construcciones de \texttt{String} que
terminen pasadas a un método de ejecución de SQL, distinguiendo
explícitamente -- como ya advirtió el
\autoref{cap:implementacion} -- entre concatenación real de variables
Java dentro del literal de la consulta (el patrón a marcar) y
\texttt{@Query} nativa con parámetros nombrados bindeados (el patrón ya
verificado como seguro) -- un análisis ingenuo basado solo en buscar la
cadena \texttt{nativeQuery = true}, sin esa distinción, marcaría este
patrón como falso positivo.

\section{Auditoría ZAP baseline scan}
\label{sec:tf-zap}

El \autoref{cap:resultados} (\autoref{sec:res-seguridad}) declaró la
auditoría automatizada OWASP ZAP como \textbf{aún no ejecutada},
asignada a Cajas según el plan del equipo. Quedó fuera de alcance de esta
entrega porque las 6 verificaciones manuales control-por-control ya
cerradas (\autoref{tab:res-owasp}) se priorizaron primero, al ser más
directamente accionables sobre hallazgos concretos que un escaneo
automatizado de superficie amplia. La línea de trabajo futura es ejecutar
\texttt{zap-baseline.py} (o el contenedor oficial
\texttt{owasp/zap2docker-stable}) contra el stack Docker real en un
entorno no productivo, documentar cada alerta con el mismo criterio de
honestidad ya aplicado al resto de este documento (real vs.\ falso
positivo, corregido vs.\ diferido), y consolidar su resultado junto al de
las 6 verificaciones manuales ya existentes -- sin descartar alertas sin
justificación documentada.

\section{Completar Lighthouse a 6 corridas (3 móvil + 3 escritorio)}
\label{sec:tf-lighthouse}

El Bloque~4 de evaluación empírica (\autoref{sec:res-lighthouse}) declaró
explícitamente que solo existen 2 corridas totales, ambas de perfil
móvil, frente a las 3 corridas por perfil (móvil + escritorio, 6 en
total) que exige la guía -- una limitación reconocida, no un resultado
completo presentado como si lo fuera. Quedó fuera de alcance de esta
entrega porque las 2 corridas existentes ya bastaron para identificar y
corregir el único hallazgo real que tenían (SEO $82\to100$,
\autoref{sec:res-lighthouse}), y ampliar la muestra no cambia una
corrección ya aplicada. La línea de trabajo futura es ejecutar 2 corridas
adicionales de perfil móvil (para completar 3) y 3 corridas de perfil
escritorio (\texttt{formFactor: desktop}, sin \emph{throttling} de Slow~4G,
en \texttt{frontend-angular/lighthouserc.js}, ya preparado para un
\emph{runner} Linux), y entonces sí calcular media/desviación
típica/IC~95\% por perfil -- una estadística que, con $n=2$ en un único
perfil, este documento decidió no forzar (\autoref{sec:res-lighthouse}).

\section{Defensa CSRF dedicada para \texttt{/api/auth/refresh}}
\label{sec:tf-csrf}

El ADR-007 (\texttt{docs/adr/adr-007-cookies-jwt.md}) documenta que la
cookie del \texttt{refreshToken} pasó de \texttt{SameSite=Strict} a
\texttt{SameSite=None} para que el navegador la enviara en peticiones
cross-origin entre el frontend
(\texttt{biblora-sgb.onrender.com}) y el backend en Render, requisito
para que el \emph{silent refresh} funcionara en producción. Verificado
en el código actual (\texttt{SecurityConfig.java}, línea 55), la
protección CSRF nativa de Spring Security está desactivada
(\texttt{.csrf(csrf -> csrf.disable())}), y no existe en el
repositorio ningún mecanismo de token CSRF ni validación manual de
\texttt{Origin}/\texttt{Referer} (verificado por búsqueda directa, cero
resultados). El único control activo es la lista blanca de CORS, que
impide que un origen no autorizado \emph{lea} la respuesta de
\texttt{/api/auth/refresh}, pero no impide que dicho origen
\emph{envíe} la petición con la cookie adjunta -- el navegador la
adjunta igual por ser \texttt{SameSite=None}. Quedó fuera de alcance de
esta entrega porque corregirlo exige un cambio coordinado
backend/frontend (emitir el token en un lugar legible por JavaScript,
p.\,ej. otra cookie no \texttt{HttpOnly} o un header de respuesta, y
que \texttt{jwt.interceptor.ts} lo reenvíe en cada petición de refresh)
que no se priorizó frente a los hallazgos de mayor severidad ya
cerrados en esta entrega (ADR-007 original, ADR-003). Se documenta
explícitamente como riesgo aceptado, no como algo ya mitigado por el
diseño actual (ver ADR-007, sección \emph{Consequences}). La línea de
trabajo futura concreta es implementar un token CSRF de patrón
\emph{double-submit cookie}: al emitir el \texttt{refreshToken}, emitir
también un valor aleatorio en una segunda cookie legible por
JavaScript (\texttt{SameSite=None} pero sin \texttt{HttpOnly}); el
frontend lo reenvía en un header custom (p.\,ej.
\texttt{X-CSRF-Token}) en la petición de refresh, y el backend rechaza
la petición si el valor del header no coincide con el de la cookie --
un atacante cross-site puede lograr que el navegador adjunte la cookie,
pero no puede leer su valor para construir el header a mano.

\section{Salvaguarda contra auto-degradación del rol ADMIN}
\label{sec:tf-admin-safeguard}

Al revisar el código real de \texttt{UsuarioAdminService.cambiarRol}
(\texttt{backend-springboot/src/main/java/com/uteq/backend/service/UsuarioAdminService.java},
método \texttt{cambiarRol}, líneas 79--95) para este capítulo se
confirmó un gap de seguridad real: el método reemplaza el conjunto de
roles de \emph{cualquier} \texttt{usuarioId} recibido por parámetro, sin
comparar ese id contra el del ejecutor autenticado
(\texttt{Authentication}) ni verificar si es el último usuario con rol
\texttt{ADMIN} del sistema. En la práctica, un ADMIN autenticado puede
degradarse a sí mismo -- o degradar al único otro ADMIN existente -- sin
ninguna salvaguarda, quedando el sistema potencialmente sin ningún
usuario con ese rol y sin una vía de recuperación dentro de la propia
aplicación. Quedó fuera de alcance de corregirse en esta entrega porque
se detectó durante la revisión de este mismo capítulo, no antes -- se
documenta aquí en vez de corregirse apresuradamente sin las pruebas
correspondientes, siguiendo el mismo criterio ya aplicado en otros
hallazgos de este documento (declarar antes que parchear sin verificar).
La línea de trabajo futura es agregar, en \texttt{cambiarRol}, una
comprobación explícita de que (a) el usuario objetivo no sea el mismo
ejecutor autenticado cuando el nuevo rol no incluya \texttt{ADMIN}, y (b)
si lo es, que exista al menos otro usuario activo con rol \texttt{ADMIN}
antes de aplicar el cambio -- junto con un test de seguridad dedicado
(análogo a \texttt{UsuarioAdminServiceTest} ya existente) que reproduzca
ambos casos y falle si la salvaguarda se elimina en el futuro.
